A mobile GUI agent often has to type text it saw many screens earlier: it reads ``vanilla extract'' on a recipe page and must enter it into a shopping list nine steps later, after that screen has left the display. We call such a step a \emph{delayed information need}. The hard part is not storing the history but deciding which earlier screen to re-read, and training a model to decide it normally requires human labels. We show that no labels need to be bought. The automatic verifier used to \emph{mine} these cases (the typed text appears in the OCR of an earlier screen and not on the current one) also \emph{labels} every earlier screen as holding the text or not, at the cost of one substring test per screen. Training on those labels, which we call \emph{verifier-supervised look-back}, is what makes look-back work. It lifts a full-history reader from 54.8\% to 65.4\% source-screen localization on 188 verified cases and a per-screen scorer from 50.0\% to 72.3\%. The two trained variants are statistically indistinguishable (\pval{0.086} on localization, \pval{0.45} end-to-end), so the supervision rather than the retrieval architecture is the lever; the per-screen decomposition pays only at gaps of 8--12 screens (78.2\% vs.\ 56.4\%, \pval{0.003}). On a fourth sample of trajectories mined after every design decision was fixed and evaluated once, the two reach 79.6\% and 76.3\% against 59.1\% for prompting (\pval{0.003}), and raise end-to-end accuracy from 23.7\% to 49.5\% and 51.6\% ($p<0.001$). We also report where both trained variants fail, and two evaluation filters without which this problem cannot be measured at all.

\textbf{Code, data artifacts and all result files:} \url{https://github.com/WMSlalalala/verifier-supervised-lookback}
